% Full text of statutes and regulations noticed in
% RJN-SUPPLEMENTAL-42517-CROSS-HEARING-{801,802}.
% Reproduced verbatim to satisfy Cal. Rules of Court, rule 3.1306(c).
% Item numbering matches the RJN's enumerated list.

\subsubsection*{Item 1. Code of Civil Procedure \S~425.17 (full text)}

\begin{quoting}
\noindent (a) The Legislature finds and declares that there has been a disturbing abuse of Section 425.16, the California Anti-SLAPP Law, which has undermined the exercise of the constitutional rights of freedom of speech and petition for the redress of grievances, contrary to the purpose and intent of Section 425.16. The Legislature finds and declares that it is in the public interest to encourage continued participation in matters of public significance, and that this participation should not be chilled through abuse of the judicial process or Section 425.16.

\noindent (b) Section 425.16 does not apply to any action brought solely in the public interest or on behalf of the general public if all of the following conditions exist:

\noindent (1) The plaintiff does not seek any relief greater than or different from the relief sought for the general public or a class of which the plaintiff is a member. A claim for attorney's fees, costs, or penalties does not constitute greater or different relief for purposes of this subdivision.

\noindent (2) The action, if successful, would enforce an important right affecting the public interest, and would confer a significant benefit, whether pecuniary or nonpecuniary, on the general public or a large class of persons.

\noindent (3) Private enforcement is necessary and places a disproportionate financial burden on the plaintiff in relation to the plaintiff's stake in the matter.

\noindent (c) Section 425.16 does not apply to any cause of action brought against a person primarily engaged in the business of selling or leasing goods or services, including, but not limited to, insurance, securities, or financial instruments, arising from any statement or conduct by that person if both of the following conditions exist:

\noindent (1) The statement or conduct consists of representations of fact about that person's or a business competitor's business operations, goods, or services, that is made for the purpose of obtaining approval for, promoting, or securing sales or leases of, or commercial transactions in, the person's goods or services, or the statement or conduct was made in the course of delivering the person's goods or services.

\noindent (2) The intended audience is an actual or potential buyer or customer, or a person likely to repeat the statement to, or otherwise influence, an actual or potential buyer or customer, or the statement or conduct arose out of or within the context of a regulatory approval process, proceeding, or investigation, except where the statement or conduct was made by a telephone corporation in the course of a proceeding before the California Public Utilities Commission and is the subject of a lawsuit brought by a competitor, notwithstanding that the conduct or statement concerns an important public issue.

\noindent (d) Subdivisions (b) and (c) do not apply to any of the following:

\noindent (1) Any person enumerated in subdivision (b) of Section 2 of Article I of the California Constitution or Section 1070 of the Evidence Code, or any person engaged in the dissemination of ideas or expression in any book or academic journal, while engaged in the gathering, receiving, or processing of information for communication to the public.

\noindent (2) Any action against any person or entity based upon the creation, dissemination, exhibition, advertisement, or other similar promotion of any dramatic, literary, musical, political, or artistic work, including, but not limited to, a motion picture or television program, or an article published in a newspaper or magazine of general circulation.

\noindent (3) Any nonprofit organization that receives more than 50 percent of its annual revenues from federal, state, or local government grants, awards, programs, or reimbursements for services rendered.

\noindent (e) If any trial court denies a special motion to strike on the grounds that the action or cause of action is exempt pursuant to this section, the appeal provisions in subdivision (i) of Section 425.16 and paragraph (13) of subdivision (a) of Section 904.1 do not apply to that action or cause of action.
\end{quoting}

\subsubsection*{Item 2. Code of Civil Procedure \S~425.16(e), (g), (i) (operative subdivisions)}

\begin{quoting}
\noindent (e) As used in this section, ``act in furtherance of a person's right of petition or free speech under the United States Constitution or the California Constitution in connection with a public issue'' includes: (1) any written or oral statement or writing made before a legislative, executive, or judicial proceeding, or any other official proceeding authorized by law; (2) any written or oral statement or writing made in connection with an issue under consideration or review by a legislative, executive, or judicial body, or any other official proceeding authorized by law; (3) any written or oral statement or writing made in a place open to the public or a public forum in connection with an issue of public interest; or (4) any other conduct in furtherance of the exercise of the constitutional right of petition or the constitutional right of free speech in connection with a public issue or an issue of public interest.

\noindent (g) All discovery proceedings in the action shall be stayed upon the filing of a notice of motion made pursuant to this section. The stay of discovery shall remain in effect until notice of entry of the order ruling on the motion. The court, on noticed motion and for good cause shown, may order that specified discovery be conducted notwithstanding this subdivision.

\noindent (i) An order granting or denying a special motion to strike shall be appealable under Section 904.1.
\end{quoting}

\subsubsection*{Item 3. Code of Civil Procedure \S~904.1(a)(13) (interlocutory appeal list)}

\begin{quoting}
\noindent (a) An appeal, other than in a limited civil case, is to the court of appeal. An appeal, other than in a limited civil case, may be taken from any of the following: \dots

\noindent (13) From an order granting or denying a special motion to strike under Section 425.16.
\end{quoting}

\subsubsection*{Item 4. Civil Code \S~47(b) (litigation privilege)}

\begin{quoting}
\noindent A privileged publication or broadcast is one made: \dots

\noindent (b) In any (1) legislative proceeding, (2) judicial proceeding, (3) in any other official proceeding authorized by law, or (4) in the initiation or course of any other proceeding authorized by law and reviewable pursuant to Chapter 2 (commencing with Section 1084) of Title 1 of Part 3 of the Code of Civil Procedure, except as follows:

\noindent (1) An allegation or averment contained in any pleading or affidavit filed in an action for marital dissolution or legal separation made of or concerning a person by or against whom no affirmative relief is prayed in the action shall not be a privileged publication or broadcast as to the person making the allegation or averment within the meaning of this section unless the pleading is verified or affidavit sworn to, and is made without malice, by one having reasonable and probable cause for believing the truth of the allegation or averment and unless the allegation or averment is material and relevant to the issues in the action.

\noindent (2) This subdivision does not make privileged any communication made in furtherance of an act of intentional destruction or alteration of physical evidence undertaken for the purpose of depriving a party to litigation of the use of that evidence, whether or not the content of the communication is the subject of a subsequent publication or broadcast which is privileged pursuant to this section. As used in this paragraph, ``physical evidence'' means evidence specified in Section 250 of the Evidence Code or evidence that is property of any type specified in Chapter 14 (commencing with Section 2031.010) of Title 4 of Part 4 of the Code of Civil Procedure.

\noindent (3) This subdivision does not make privileged any communication made in a judicial proceeding knowingly concealing the existence of an insurance policy or policies.

\noindent (4) A recorded lis pendens is not a privileged publication unless it identifies an action previously filed with a court of competent jurisdiction which affects the title or right of possession of real property, as authorized or required by law.
\end{quoting}

\subsubsection*{Item 5. Business and Professions Code \S~6128(a) (attorney deceit)}

\begin{quoting}
\noindent Every attorney is guilty of a misdemeanor who either:

\noindent (a) Is guilty of any deceit or collusion, or consents to any deceit or collusion, with intent to deceive the court or any party.
\end{quoting}

\subsubsection*{Item 6. Insurance Code \S~790.03(h)(1), (h)(3), (h)(5) (unfair settlement practices)}

\begin{quoting}
\noindent The following are hereby defined as unfair methods of competition and unfair and deceptive acts or practices in the business of insurance. \dots

\noindent (h) Knowingly committing or performing with such frequency as to indicate a general business practice any of the following unfair claims settlement practices:

\noindent (1) Misrepresenting to claimants pertinent facts or insurance policy provisions relating to any coverages at issue. \dots

\noindent (3) Failing to adopt and implement reasonable standards for the prompt investigation and processing of claims arising under insurance policies. \dots

\noindent (5) Not attempting in good faith to effectuate prompt, fair, and equitable settlements of claims in which liability has become reasonably clear.
\end{quoting}

\subsubsection*{Item 7. Title 10, California Code of Regulations, \S~2695.1 (Fair Claims Settlement Practices Regulations --- scope and purpose)}

\begin{quoting}
\noindent (a) Section 790.03(h) of the California Insurance Code enumerates sixteen claims settlement practices that, when either knowingly committed on a single occasion, or performed with such frequency as to indicate a general business practice, are considered to be unfair claims settlement practices and are, thus, prohibited by this section of the California Insurance Code. The regulations which follow set forth the minimum standards for the settlement of claims and shall be applicable to the handling or settlement of all claims, except as specifically excluded.

\noindent (b) These regulations are intended to delineate certain minimum standards for the settlement of claims which, when violated knowingly on a single occasion or performed with such frequency as to indicate a general business practice shall constitute an unfair claims settlement practice within the meaning of Insurance Code Section 790.03(h). It is also the intent of these regulations to promote the good faith, prompt, efficient and equitable settlement of claims on a cost-effective basis.

\noindent (c) These regulations are not intended to provide the basis for an institutional general business practice violation of California Insurance Code Section 790.03(h).

\noindent (d) These regulations are applicable to all persons and to all insurance policies, insurance contracts, and certificates, except as specifically excluded.

\noindent The full Fair Claims Settlement Practices Regulations compilation (10 CCR \S\S~2695.1--2695.17) is concurrently lodged with this Request.
\end{quoting}

\subsubsection*{Item 8. Insurance Code \S~790.02 (legislative purpose --- excerpt)}

\begin{quoting}
\noindent The purpose of this article is to regulate trade practices in the business of insurance in accordance with the intent of Congress as expressed in the Act of Congress of March 9, 1945 (Public Law 15, Seventy-ninth Congress), by defining, or providing for the determination of, all such practices in this State which constitute unfair methods of competition or unfair or deceptive acts or practices and by prohibiting the trade practices so defined or determined.
\end{quoting}

\subsubsection*{Item 9. Title 10, California Code of Regulations, \S\S~2695.4(a), 2695.5, 2695.7(a), (b), and (d) (excerpts)}

\begin{quoting}
\noindent \textbf{2695.4(a).} Every insurer shall disclose to a first party claimant or beneficiary, all benefits, coverage, time limits or other provisions of any insurance policy issued by that insurer that may apply to the claim presented by the claimant. When additional benefits might reasonably be payable under an insured's policy upon receipt of additional proofs of claim, the insurer shall immediately communicate this fact to the insured and cooperate with and assist the insured in determining the extent of the insurer's additional liability.

\noindent \textbf{2695.5.} Upon receiving notification of a claim, every insurer shall immediately provide the claimant necessary claim forms, instructions, including reasonable assistance so that the claimant can make claims properly with respect to all coverages and matters related to the claim.

\noindent \textbf{2695.7(a).} No insurer shall discriminate in its claims settlement practices based upon the claimant's age, race, gender, income, religion, language, sexual orientation, ancestry, national origin, or physical disability, or upon the territory of the property or person insured.

\noindent \textbf{2695.7(b).} Upon receiving proof of claim, every insurer, except as specified in subsection 2695.7(b)(4) below, shall immediately, but in no event more than forty (40) calendar days later, accept or deny the claim, in whole or in part. \dots

\noindent \textbf{2695.7(d).} Every insurer shall conduct and diligently pursue a thorough, fair and objective investigation and shall not persist in seeking information not reasonably required for or material to the resolution of a claim dispute.
\end{quoting}

\subsubsection*{Item 10. Civil Code \S~2330 (agency)}

\begin{quoting}
\noindent An agent represents his principal for all purposes within the scope of his actual or ostensible authority, and all the rights and liabilities which would accrue to the agent from transactions within such limit, if they had been entered into on his own account, accrue to the principal.
\end{quoting}
